\subsection{How to read the output}\label{sm:reading}
The worked example of \cref{sec:worked}, run with the benchmark protocol (one noisy grid, $16$ bootstrap refits),
produces a text report and five images; all of them are in the folder \texttt{example\_run/} of the archive, and the
report returns the equation of \cref{tab:worked}. This section says how to read them. With scattered data the same
report is computed over posterior samples of the field instead of bootstrap refits.

\paragraph{Three kinds of number} The output mixes three kinds of number, which must not be confused.
\emph{Posterior probabilities} ($P(k)$, $P(\Mfit)$, $P(\text{split})$) come from model comparison and say which model wins
within the models considered; they are not p-values. \emph{Goodness-of-fit p-values}, always printed as $\log_{10}p$, test
the null hypothesis that the reported equation holds and the residual is white noise of the given level; a small $p$
withdraws the claim (below $10^{-3}$ the label becomes \MIS), while a large $p$ only says that no misfit was found, which
is why the label is \FC{} and never ``identified''. \emph{Distributions over posterior samples} (of $\cF$, $\gN$, $\rho$,
the shares $w_j$ and every coefficient) are summarized by their $5\%$, $50\%$ and $95\%$ quantiles; read the width of the
interval first.

\paragraph{Directions} A factor $Q$ can carry the mode $e^{\lambda t+i\kappa x}$ exactly when $Q(\lambda,i\kappa)=0$. The image
\texttt{directions.png} (\cref{sm:fig:directions}) therefore shows, for each factor, $\log_{10}\abs{Q(\lambda,i\kappa)}$ over the
plane of wavenumber and rate: dark pixels are the modes the factor can carry, and reading a factor means reading its
dark curve. Because $\lambda$ is complex each factor gets two panels, oscillation ($\operatorname{Im}\lambda$) and damping
($\operatorname{Re}\lambda$). A tilted line in the oscillation panel with a flat line at zero in the damping panel is
transport; a flat line in the oscillation panel with a downward parabola in the damping panel is diffusion; curves in
both are damped waves. A blurred or broken curve means poorly determined coefficients.

\begin{figure}[htbp]
\centering
\includegraphics[width=0.8\linewidth]{sm_directions}
\caption{\texttt{directions.png}: each column is one reported factor; dark curves are the modes it can carry. Left:
transport (tilted in the oscillation row, flat in the damping row). Right: diffusion (flat in the oscillation row, a
parabola in the damping row).}\label{sm:fig:directions}
\end{figure}

\paragraph{Separation of the factors} The image \texttt{operator\_space.png} (\cref{sm:fig:opspace}) compares the reported
factors and the factors found in the posterior samples by the largest principal angle between their solution spaces. Two
compact blocks far apart mean two well separated factors on which the samples agree; a spread-out cloud, or samples
halfway between two factors, means that the samples disagree about the structure; two reported factors close together
mean a small $\gN$ and wide intervals for the components.

\begin{figure}[htbp]
\centering
\includegraphics[width=\linewidth]{sm_operator_space}
\caption{\texttt{operator\_space.png}. Left: distance between operators. Middle: largest principal angle between the
solution spaces of every pair of factors (the reported factors \texttt{f0}, \texttt{f1} and the factors \texttt{s0t0},
\texttt{s0t1} found in the field sample), ordered by cluster; dark diagonal blocks are the groups, about $1.2$ rad apart.
Right: a two-dimensional map of these angles; stars are the reported factors, dots the factors of the
samples.}\label{sm:fig:opspace}
\end{figure}

\paragraph{Components} The image \texttt{modes\_and\_components.png} shows, for each factor, a few of its kernel modes and
the component $u_j$ of the data attributed to it; the components add up to the data up to the residual. If a component
looks like noise, or like a copy of another, that factor is not well separated. \texttt{directions\_uncertainty.png}
shows where the posterior samples disagree about each curve, and \texttt{distributions.png} the histograms of $\cF$, $\gN$
and $\rho$ over the samples.

\paragraph{Five questions to ask of any result}
\begin{enumerate}
\item What is the label? Anything other than \FC: read the note that comes with it.
\item Are the factors separated? $P(\text{split})$ close to one, $\gN$ not tiny, two distinct blocks in
\cref{sm:fig:opspace}.
\item Is each coefficient narrow and nonzero? Look at the $90\%$ interval of every term.
\item Does the model pass the tests? No $\log_{10}p$ below $-3$.
\item Do the pictures look like physics? Clean dark curves, components that look like the data, posterior samples on top
of the reported factors.
\end{enumerate}

